\chapter{Review of topology in $\mathbb{R}^n$}


\section{Compactness in $\mathbb{R}^n$}
\subsection{Finite Intersecion property and Heine-Borel}
\begin{theorem}
    \textbf{cpt 集的任意闭子集都继承 cpt 性.}
\end{theorem}

\begin{theorem}{Finite Intersecion property} \label{Finite Intersecion property}
   任意一堆(even unctbl个) cpt 集，只要知道它们中\textbf{任意有限个的交集都不是空的，那么它们整体的交集也不是空的. }
\end{theorem}


\begin{theorem}
    \textbf{Metric space 中，cpt 和 seq cpt 是等价的.}
    （但是一般的 topological space 中，并不等价且无蕴含关系）
    review: sequential compact 就是这个集合中的任意序列一定存在收敛子序列，并且收敛至这个集合中的一点.
\end{theorem}

\begin{theorem}{Nested Box property} \label{Nested Box property}
    这是 $\mathbb{R}^n$ 的一个性质。意思是 $\mathbb{R}^n$ 中的任意 nested 的闭区间的序列，$\cap_{i=1} ^\infty [a_i, b_i]$ 一定是非空的. (可能是一点，也可能是一个闭区间.) 
\end{theorem}

\begin{theorem}{Heine-Borel} \label{Heine-Borel}
    textbf{$\mathbb{R}^n$ 中，closed + bdd 等价于 cpt (等价于 seq cpt).} \\
（但是这个性质非常局限，因为它甚至不能推广至 metric space. 它在普遍的 metric space 中并不成立.）
\end{theorem}

一个重要例子：
\subsection{$l^{\infty}(\mathbb{N})$}

\begin{definition}{$l^{\infty}(\mathbb{N})$} \label{linfinity}
    所有 bouded sequence in $\mathbb{R}$ 构成的空间，其中 $\mathbb{N}$ 的意思是 index on自然数集，也可以取其他 countable 的集合来作为 index: $l^{\infty}(X)$ .
\end{definition}


这个空间采取\textbf{ $\sup$ norm metric：$d((x_n), (y_n)) = \sup_{n} |x_n - y_n| $. 也就是这两个 sequence 逐项差的 sup.}

这个 metric space 中的单位球是 closed 且 bounded 的(显然).  

BTW: 这个 $l^{\infty}(\mathbb{N})$ 空间也是一个 infinite dimension 的 vector space，我们在 LADR 中已经学过. 并且，我们可以给他赋范，也就是 sup norm 的由来：\textbf{$||x_n||_{\infty} = sup_{n}|x_n|$}，而它的 metric 其实是它的 sup norm induce 出来的. 这个 sup norm 其实就是这个序列和零序列之间的 sup norm metric. 

\textbf{并且，在 sup norm 下作为一个 normed vector space，它是一个 Banach space。
（hw1: 任意的 normed vector space 在 $d =  ||x-y||$ 下都是一个 metric space. 如果这个 norm 使得它是 complete metric space，那么就称之为一个 Banach space.）}

然而，这个单位球并不是 cpt 的。我们知道，metric space 中 seq cpt 等价于 cpt. 所以我们可以构造一个序列 in $B_1(0)$: 考虑 $((1,0,0,\cdots), (0,1,0,\cdots),(0,0,1,\cdots),\cdots)$ ，其中每个序列的第 $n$ 项是 1 其他都是 0. 很显然这个序列并没有收敛子序列，因为其中每一项和其他之间 $d_{\infty}$ 都是1. 所以它并不 cpt.


\section{Totally-Compactness and Completeness}
\subsection{cpt. = complte + ttlbdd.}
我们知道在 $\mathbb{R}^n$ 里面，cpt = seq cpt = closed + bounded.
但是\textbf{在一般的 metric space 里面，虽然还是有 cpt = seq cpt，却不等价于有界闭了。}

不过我们仍然希望把 cptness 这个 topology 的概念转换成更加直观的概念以服务于 analysis 的目的。Fortunately 这是可行的：我们现在定义

\begin{definition}{Totally Bounded} \label{ttlbdd}
    任取 $\epsilon > 0$，我们都可以用 finitely many 个 $\epsilon$-ball 来 cover 这个集合，则称它是 Totally Bounded 的.
\end{definition}

\begin{definition}{Complete (或称 Cauchy-complete)} \label{cplt}
    即这个集合中的任意 Cauchy sequence 都收敛（至这个集合里的某个点）
    \\Rmk: 我们知道 \textbf{metric space 的任意子集也是 metric space}，by induced metric.
\end{definition}

\begin{remark}
complete 即这个集合里任意 Cauchy 序列都收敛于它自己的某个点，ttl bounded 即对于任意半径都可以用有限个这样半径的球来覆盖这个集合.
\end{remark}



现在我们证明这个重要定理：
\begin{theorem}
    \textbf{任意 metric space 中，cpt = seq cpt = complete + ttl bounded}
\end{theorem}

\subsection{Lebesgue Covering Thm}
为了证明这个东西我们需要一个 Lemma，其实这个 Lemma 叫做 \textbf{Lebesgue Covering Thm}:  
\begin{lemma}{Lebesgue Covering Thm} \label{lbg covering thm}
    对于 cpt 的 $E$ ，取任意一个 open cover $\{U_{\alpha}\}$，都存在一个 $\epsilon > 0$ 使得对于 $E$ 中的任何一点 $p$，$B_{\epsilon}(p)$ 都被包含在某个 $U_{\alpha}$ 中.
\end{lemma}
\proof 这个定理其实非常直观。因为 compact 一定有 finite subcover. 既然是 finite 的，那么任意一个 cover set 和集合或其他 cover set 边缘擦边的地方一定是有限的。既然 $E$ 是 compact 的那么一定是 complete 的那么当然也是 closed 的，于是取这些 finite 个擦边的地方里面最擦边的地方作为 $\epsilon$ 就可以了.

\begin{remark}
    显然，metric space 中 \textbf{complete 蕴含 closed，ttl bdd 蕴含 bdd}，而特别地， $\mathbb{R}^n$ 中的任意集合，closed 等价于 complete，ttl bdd 等价于 bdd.
\end{remark}

\begin{theorem}
    \textbf{complete 的 metric space 的任意 closed subset 一定也是 complete 的}. 显然.
\end{theorem}
complete 的 metric space 的任意 closed subset（我们知道 metric space 的任意子集也是 metric space，by induced metric.）一定也是 complete 的. 显然.

\begin{remark}
    我们在 lec 1 中讲到  $l^{\infty}(\mathbb{N})$ 中的单位球是 closed + bdd 但是不 compact. 由于这是一个 Banach space 它也是 complete 的，by 3；不过它的问题出在并不是 ttl bdd 的.
\end{remark}


